\chapter{Introduction}
\label{chap:intro}

\section{Context}
\label{sec:context}

Environmental surveillance beacons monitor radiological pollution by counting,
at a fixed sampling interval (one minute here), the gamma photons that reach a
scintillation detector. The resulting \emph{gamma-dose} time series reflects the
natural radioactive background, modulated by weather and, rarely, by a
radiological event: the passage of a contaminated source, a leak, or the
dispersal of a radioactive plume. The operational goal is to detect and
\emph{identify} such events as early as possible, directly from the shape they
leave in the signal. Three properties make this harder than textbook anomaly
detection: the per-sample noise is comparable in size to weak events; the
baseline drifts slowly with meteorological conditions, so no fixed threshold
works; and different physical phenomena leave \emph{different} temporal
signatures, so detection alone is not enough --- the response must classify the
shape.

\section{Problem formalisation}
\label{sec:problem}

We treat the problem as supervised classification on a univariate signal.
Let $x = (x_1,\dots,x_N) \in \R^{N}$ be the gamma-dose series sampled at uniform
intervals. Each instant $t$ belongs to one of $K$ classes: a \texttt{Normal
Background} class (label $0$) and $K-1$ anomaly morphologies. The instantaneous
value $x_t$ carries almost no class information --- a class is defined by the
local \emph{shape} of the signal --- so we classify fixed-length windows rather
than samples.

\paragraph{Windowing.} For a window length $W$ we map each position $t$ to the
sub-sequence
\begin{equation}
    \mathbf{x}^{(t)} = (x_{t-W+1},\,\dots,\,x_t) \in \R^{W},
    \label{eq:window}
\end{equation}
and ask a model $f_\theta:\R^{W}\!\to\R^{K}$, with parameters $\theta$, to map it
to class logits whose softmax
$\hat{p}^{(t)} = \softmax\!\big(f_\theta(\mathbf{x}^{(t)})\big) \in \Delta^{K-1}$
is a distribution over the $K$ classes.\footnote{$\Delta^{K-1}$ is the
probability simplex, i.e.\ the set of non-negative vectors of $\R^{K}$ summing to
$1$.} We use $W=512$ samples ($\approx 8.5$ hours). Each window receives a single
label $y^{(t)}\in\{0,\dots,K-1\}$ through the coverage rule of
\cref{sec:windowing}.

\paragraph{Learning objective.} Given a training set
$\mathcal{D}=\{(\mathbf{x}_i,y_i)\}$ of labelled windows, $\theta$ is fitted by
minimising the class-weighted cross-entropy
\begin{equation}
    \mathcal{L}(\theta)
    = -\frac{1}{|\mathcal{D}|}\sum_{(\mathbf{x},y)\in\mathcal{D}}
        w_y \,\log \softmax\!\big(f_\theta(\mathbf{x})\big)_y ,
    \qquad
    w_c = \frac{|\mathcal{D}|}{K\,n_c},
    \label{eq:objective}
\end{equation}
where $n_c$ is the number of training windows of class $c$ and the weight $w_c$
compensates the strong class imbalance (\cref{sec:training}).

\paragraph{Evaluation objective.} Sample- or window-level accuracy over-counts
the many background windows and the fuzzy boundaries of an event. The
operationally meaningful quantity is whether each \emph{event} is detected and
correctly named. We therefore also report an \emph{event-level} macro
$F_1$-score (\cref{sec:event-eval}): predictions are aggregated into a per-sample
confidence trace, contiguous non-background runs form predicted events, and each
is matched against the ground-truth events. This bridges window classification
and the detection task the beacon actually performs.

\section{Why deep learning, and what we validate}
\label{sec:why-dl}

A classical pipeline --- baseline removal followed by a bank of hand-tuned
matched filters --- works on clean, isolated shapes but scales poorly: real
events are time-warped and noisy, cross-correlation with a template degrades
quickly at the observed noise level, and every new morphology requires
redesigning the bank. Convolutional and attention-based networks instead
\emph{learn} their own shape features from data and are, by now, a strong
baseline for time-series classification (\cref{chap:related}). The project brief
\cite{centralesupelec_projet} explicitly suggests evaluating convolutional and
attention layers, which is what we do.

Crucially, \textbf{no labelled real events are available}. Following the
methodological choice that a deliberately simplified, well-understood synthetic
model is the right tool to \emph{validate an approach}, we generate a synthetic
benchmark whose noise statistics are calibrated to real recordings, and we
evaluate the networks only on held-out \emph{synthetic} data. Running networks
trained on a simplified model against real signals would measure the modelling
gap, not the classifier, and is deliberately left as future work
(\cref{chap:conclusion}).

\section{Contributions and outline}
\label{sec:contrib}

This report (i) formalises radiological-event identification as windowed
multi-class classification (\cref{sec:problem}); (ii) positions the work against
the deep-learning time-series literature (\cref{chap:related}); (iii) describes a
simplified, statistically calibrated synthetic gamma-dose generator
(\cref{chap:data}); (iv) specifies three architectures of increasing complexity
and a shared training and event-evaluation pipeline (\cref{chap:methods}); and
(v) compares them on classification performance, robustness and computational
cost (\cref{chap:results}). \Cref{chap:conclusion} concludes and lays out the
path towards more realistic signal models and real-data evaluation.
